\begin{singlespace}
{\itshape

\indent Chronic coronary disease (CCD) carries the largest single share of the global
cardiovascular burden, yet guideline adherence in routine practice remains poor: of the
4{,}548 coronary patients in the INTERASPIRE study, only 1.0\% met its definition of optimal
adherence. A recommendation system can close that gap only if it solves three problems at
once---fusing the patient record, the clinical guideline and the biomedical literature into
one queryable structure, bridging the terminologies in which each is written, and generating
natural-language output whose faithfulness can actually be checked against the source
record. This thesis develops a FHIR-native, knowledge-graph-based recommendation system
for CCD that addresses all three.

Synthetic records for 5{,}490 patients were generated with Synthea, exported as HL7 FHIR~R4
bundles, loaded into a Neo4j property graph, and fused with the 2023 AHA/ACC CCD guideline
and 4{,}801 PubMed abstracts through a UMLS concept layer. Because numeric observations are
loaded as graph data rather than discarded, recommendations conditioned on measured
values---ejection fraction or body-mass index, for instance---are gated on what the patient
was actually measured at, which a condition-only pipeline cannot do. The system also
produces a grounded briefing in which every fact is retrieved deterministically from the
graph and the language model does nothing but phrase it, after which an automated check
verifies the result. Evaluation covered 876 CCD patients, scored against a reference
standard generated independently of the system.

The resulting graph holds 1{,}804{,}379 nodes and 8{,}796{,}425 relationships. An ablation
shows that observation gating raises guideline concordance $F_1$ from 0.743 to 0.810,
through precision rising from 0.612 to 0.684 with recall essentially unchanged, and the
whole gain is confined to the 15 of 68 recommendations that are gated at all. The automated
reference standard agrees with manual labelling of 30 patients on 99.5\% of cells, Cohen's
$\kappa = 0.991$ (95\% CI 0.984--0.996, resampling patients). All 30 generated briefings passed the automated faithfulness check, with no
fabricated recommendation or citation, no omitted required item, and no anticoagulation
over-claim. A cardiologist judged 77.7\% of 247 recommendations for 15 patients appropriate,
with the judgements of inappropriateness concentrated on antithrombotic therapy.

Loading numeric observations as first-class data therefore raises the precision of guideline
concordance measurably. FHIR supplies those observations most directly, but a tabular export
reproduces the same gate. Deterministic retrieval makes the citation-level faithfulness of
generated clinical text decidable rather than estimated, but it does not guarantee that every
sentence is semantically faithful.
The limits are measured too: the FHIR export carries no medication stop dates, so
duration-dependent guideline logic can only be partly recovered, and every figure above
comes from synthetic records.

}
\end{singlespace}

\vspace{6pt}
\noindent\textbf{Keywords} : clinical recommendation system, knowledge graph, HL7 FHIR,
chronic coronary disease, guideline concordance
